\section{Tesis 2}

\subsection{Título}
\textit{Adaptive Intelligent User Interfaces With Emotion Recognition} (Interfaces de Usuario Inteligentes y Adaptativas con Reconocimiento de Emociones) \cite{nasoz2004}.

\subsection{Autor}
Fatma Nasoz, University of Central Florida.

\subsection{Resumen}
``El enfoque de esta disertación es la creación de Interfaces de Usuario Inteligentes y Adaptativas para facilitar una comunicación natural y enriquecida durante la Interacción Humano-Computadora, reconociendo los estados afectivos de los usuarios (es decir, las emociones experimentadas por ellos) y respondiendo a esas emociones mediante la adaptación a la situación actual a través de un modelo de usuario afectivo creado para cada usuario. Se diseñaron y realizaron experimentos controlados en un entorno de laboratorio y en un entorno de Realidad Virtual para recolectar señales de datos fisiológicos de participantes que experimentaban emociones específicas. Se implementaron algoritmos (k-Vecinos más Cercanos [KNN], Análisis de Función Discriminante [DFA], Retropropagación de Marquardt [MBP] y Retropropagación Resiliente [RBP]) para analizar las señales de datos recolectadas y encontrar patrones fisiológicos únicos de las emociones. Se realizó un Experimento de Elicitación de Emociones con Clips de Películas para elicitar Tristeza, Enojo, Sorpresa, Miedo, Frustración y Diversión de los participantes. En general, los tres algoritmos KNN, DFA y MBP pudieron reconocer emociones con una precisión del 72.3\%, 75.0\% y 84.1\%, respectivamente. Se realizó un experimento en un Simulador de Conducción para elicitar emociones relacionadas con la conducción (pánico/miedo, frustración/enojo y aburrimiento/somnolencia). Los algoritmos KNN, MBP y RBP se utilizaron para clasificar las señales fisiológicas por emociones correspondientes. En general, KNN pudo clasificar estas tres emociones con un 66.3\%, MBP con un 76.7\% y RBP con un 91.9\% de precisión. La adaptación de la interfaz se diseñó para proporcionar retroalimentación multimodal a los usuarios sobre su estado afectivo actual y para responder a los estados emocionales negativos de los usuarios con el fin de disminuir los posibles impactos negativos de esas emociones. Se empleó la formalización de Redes de Creencias Bayesianas para desarrollar el Modelo de Usuario que permite al sistema inteligente adaptarse apropiadamente al contexto y situación actuales considerando factores dependientes del usuario, como rasgos de personalidad y preferencias.'' \cite{nasoz2004}

\subsection{Problema que busca resolver}
La tesis aborda la brecha fundamental que existe en la Interacción Humano-Computadora (HCI): los sistemas informáticos son incapaces de percibir y responder a los estados emocionales de sus usuarios. Los seres humanos son afectados continuamente por sus emociones, las cuales influyen en procesos cognitivos clave como el aprendizaje, la toma de decisiones, la motivación y el desempeño. Sin embargo, las interfaces digitales convencionales ignoran completamente esta dimensión afectiva.

El problema se manifiesta con especial gravedad en tres dominios de aplicación: (1) la seguridad vial, donde emociones como el enojo, la frustración y la somnolencia son causas principales de accidentes de tránsito; (2) el aprendizaje y entrenamiento, donde la ansiedad y la frustración reducen la capacidad de aprender y la autoeficacia; y (3) la telemedicina, donde la ausencia de canales afectivos en la comunicación médico-paciente remota dificulta el diagnóstico y el tratamiento. En todos estos contextos, una interfaz que pueda reconocer las emociones del usuario y adaptarse en consecuencia mejoraría significativamente la experiencia, la seguridad y los resultados.

\subsection{Objetivos}

Objetivo General: Diseñar y construir un Sistema Inteligente Adaptativo que reconozca los estados afectivos de sus usuarios a partir de señales fisiológicas no invasivas y que, mediante un modelo de usuario, adapte la interfaz de forma personalizada para mejorar la Interacción Humano-Computadora en aplicaciones de seguridad vial, entrenamiento y telemedicina \cite{nasoz2004}.

Objetivos Específicos:
\begin{enumerate}
    \item Diseñar experimentos controlados para elicitar emociones específicas, recolectar señales fisiológicas de los participantes y registrar sus expresiones faciales; incluyendo un experimento en laboratorio con estímulos multimodales y un experimento de simulador de conducción en Realidad Virtual.
    \item Analizar las señales fisiológicas medidas implementando algoritmos de aprendizaje automático (KNN, DFA, MBP, RBP) para identificar patrones fisiológicos únicos asociados a cada emoción.
    \item Utilizar los patrones identificados para reconocer las emociones del usuario y desarrollar interfaces de usuario afectivas inteligentes que proporcionen retroalimentación apropiada sobre el estado emocional del usuario.
    \item Crear modelos de usuario e integrarlos con las interfaces desarrolladas, empleando Redes de Creencias Bayesianas, para adaptar el sistema a la emoción, la personalidad, la edad y el género del usuario en el contexto y la aplicación actuales.
\end{enumerate}

\subsection{Hipótesis}

La hipótesis no se enuncia de manera explícita en la tesis. Sin embargo, mediante ingeniería inversa a partir de los experimentos realizados, se puede inferir la siguiente hipótesis subyacente: \textit{Las señales fisiológicas del Sistema Nervioso Autónomo (respuesta galvánica de la piel, frecuencia cardíaca y temperatura cutánea), capturadas mediante dispositivos portátiles no invasivos, contienen patrones únicos para cada emoción que pueden ser identificados por algoritmos de aprendizaje automático con suficiente precisión como para alimentar un sistema adaptativo de interfaz de usuario que, al integrar un modelo de usuario basado en Redes de Creencias Bayesianas, sea capaz de elegir respuestas personalizadas apropiadas al estado emocional del usuario.}

Esta hipótesis habría sido refutada si los algoritmos de aprendizaje automático hubieran mostrado precisiones cercanas al azar (es decir, alrededor del 16.7\% para seis clases), o si los patrones fisiológicos encontrados no hubieran sido estadísticamente distinguibles entre emociones. Los resultados confirmaron la hipótesis: el mejor algoritmo (MBP) alcanzó el 84.1\% de precisión en el experimento de laboratorio y RBP alcanzó el 91.9\% en el experimento de conducción, muy por encima del nivel aleatorio \cite{nasoz2004}.

\subsection{Resumen de la metodología empleada}

La investigación empleó una metodología experimental en dos fases, apoyada en una herramienta propia del laboratorio diseñada para capturar señales corporales y adaptar la interfaz en respuesta a las emociones detectadas \cite{nasoz2004}.

\begin{itemize}
    \item \textbf{Instrumentación:} Se utilizó una pulsera portátil inalámbrica (BodyMedia SenseWear) junto con un monitor cardíaco pectoral para medir señales del cuerpo: la respuesta eléctrica de la piel (que varía con la activación emocional), la frecuencia cardíaca, la temperatura corporal, el movimiento y el flujo de calor.

    \item \textbf{Experimento 1 --- Provocación de emociones con clips de películas:} Se realizó un estudio previo con 20 participantes para seleccionar los fragmentos de video más efectivos para provocar cada emoción. Luego se recolectaron datos de participantes expuestos a seis emociones (tristeza, enojo, sorpresa, miedo, frustración y diversión) mediante una combinación de videos, imágenes, música y problemas matemáticos. Las señales se analizaron con tres algoritmos de clasificación automática.

    \item \textbf{Experimento 2 --- Simulador de conducción en Realidad Virtual:} Se diseñó un escenario de conducción virtual con eventos de tráfico que inducían tres emociones: pánico o miedo, frustración o enojo, y aburrimiento o somnolencia. Para este experimento se añadió un cuarto algoritmo de clasificación y la red neuronal se configuró con 12 variables de entrada, una capa intermedia de 17 nodos y 3 posibles salidas (una por emoción).

    \item \textbf{Modelo de usuario para personalizar la respuesta:} Se construyó un modelo probabilístico que toma como entradas la emoción detectada, la edad, el género y cinco rasgos de personalidad del usuario (neuroticismo, extraversión, apertura, amabilidad y responsabilidad), y elige la respuesta más apropiada de la interfaz para esa combinación particular.
\end{itemize}

\subsection{Principales resultados}

Los principales hallazgos reportados por \cite{nasoz2004} son:

\begin{itemize}
    \item \textbf{Experimento 1 (6 emociones en laboratorio):} Los tres algoritmos de clasificación probados lograron identificar correctamente las emociones a partir de señales corporales. El más básico alcanzó un 72.3\% de precisión, el segundo un 75.0\%, y el más sofisticado un 84.1\%. Para contexto: si el sistema eligiera al azar entre las seis emociones posibles, acertaría solo el 16.7\% de las veces, por lo que estos resultados representan una mejora considerable.

    \item \textbf{Experimento 2 (3 emociones en simulador de conducción):} En el escenario de conducción virtual, el algoritmo más efectivo alcanzó un 91.9\% de precisión global. Desglosado por emoción, identificó pánico o miedo en el 89.7\% de los casos, frustración o enojo en el 96.7\%, y aburrimiento o somnolencia en el 88.9\%.

    \item \textbf{Mejor algoritmo para el contexto de conducción:} Uno de los algoritmos de red neuronal evaluados demostró mayor capacidad de adaptarse a datos nuevos que los demás en el escenario de conducción, convirtiéndose en uno de los aportes metodológicos clave de la tesis.

    \item \textbf{Adaptación personalizada de la interfaz:} El sistema demostró poder elegir respuestas distintas según la combinación de emoción detectada, personalidad, edad y género del usuario. Por ejemplo, ante un conductor enojado mayor de 25 años, mujer y con alto sentido de responsabilidad, el sistema eligió sugerir una técnica de relajación (con 74.3\% de confianza); para un conductor enojado menor de 25 años, hombre y con tendencia a la ansiedad, la respuesta elegida fue hacer un chiste (con 58.1\% de confianza).

    \item \textbf{Sistema completo funcional:} Se implementó y demostró el sistema completo, incluyendo un avatar capaz de reflejar las emociones del usuario (felicidad, tristeza, enojo), reconocer al usuario por su cara y adaptar la interfaz a sus preferencias previamente registradas.
\end{itemize}

\subsection{Crítica de la tesis}
\begin{itemize}
    \item \textbf{Aspectos positivos:} Las secciones de la tesis están bien delimitadas y siguen una progresión lógica coherente, desde la fundamentación teórica hasta la validación experimental, lo cual facilita la lectura y comprensión del trabajo. Los objetivos específicos se enuncian de forma explícita y existe una correspondencia clara entre estos y las conclusiones: cada uno es abordado en los experimentos y retomado al cierre del trabajo. El título refleja con precisión el alcance del trabajo, que efectivamente integra el reconocimiento de emociones con la adaptación de interfaces. Asimismo, el alcance es consistente con el trabajo relacionado del área de computación afectiva e HCI, posicionando la propuesta como una contribución integradora frente a trabajos previos que solo abordaban una parte del problema.
    \item \textbf{Aspectos de posible mejora:} En primer lugar, la hipótesis no se enuncia en ninguna sección de forma explícita, lo que obliga al lector a inferirla mediante ingeniería inversa a partir de los experimentos diseñados. En segundo lugar, aunque los objetivos se alcanzan en términos generales, existe una brecha entre las conclusiones y el objetivo de adaptación: el componente basado en Redes de Creencias Bayesianas fue validado con probabilidades condicionales arbitrarias no calibradas por expertos, lo que debilita el cierre del trabajo respecto a su objetivo más distintivo. En tercer lugar, la sección denominada ``Visualización de los Resultados'' resulta confusa y mal nombrada: en lugar de presentar resultados analíticos, describe las capacidades visuales del sistema con capturas de pantalla, lo que no sigue la convención estándar de una tesis. Por último, si bien el trabajo relacionado cubre los algoritmos y sensores más relevantes, no justifica explícitamente la elección del enfoque frente a alternativas contemporáneas, dejando el alcance del trabajo parcialmente sin fundamentar en la literatura existente.
\end{itemize}
